\chapter{Deret Numerik}\label{ch:b1:series}

Menjumlahkan tak hingga banyak bilangan berarti mengambil limit jumlah
parsialnya --- tak lebih, tak kurang. Bab ini menyiapkan
definisi dan uji kekonvergenan yang terpakai di tahun pertama:
yaitu perbandingan dan kesetaraan bagi suku positif, \omterm{thm:b1:series:ratio}{uji nisbah},
perbandingan \omterm{thm:b1:integration:def}{integral} yang menghasilkan \omterm{thm:b1:series:riemann}{deret Riemann}, \omterm{thm:b1:series:absolute}{kekonvergenan mutlak},
dan teorema \omterm{thm:b1:series:alternating}{deret berselang-seling}. Adapun teori yang lebih halus (yaitu hasil kali
\omterm{def:b1:series:def}{deret}, penjumlahan berpaket, dan \omterm{def:b1:series:def}{deret} fungsi) termasuk bahan tahun
kedua.

\section{Keumuman}

\begin{definition}\label{def:b1:series:def}
Diberikan sebuah barisan $(u_n)$, \emph{deret}\index{deret} $\sum u_n$
adalah barisan \emph{jumlah parsial} $S_N = \sum_{n=0}^{N} u_n$.
Deretnya \emph{konvergen} bila $(S_N)$ konvergen; dan limitnya menjadi
\emph{jumlahnya} $\sum_{n=0}^{\infty} u_n$, sedangkan $R_N = \sum_{n > N} u_n =
S - S_N$ merupakan \emph{sisanya}, yang menuju $0$.
\end{definition}

\begin{example}[Deret geometri]\label{ex:b1:series:geometric}
Untuk $q \in \C$: $\;S_N = \sum_{n=0}^{N} q^n = \frac{1 -
q^{N+1}}{1-q}$ (dengan $q \neq 1$). \omterm{def:b1:series:def}{Deretnya} konvergen jika dan hanya jika $\abs q < 1$
(\cref{exo:b1:seq:3}), dengan\index{deret geometri}
\[
\sum_{n=0}^{\infty} q^n = \frac{1}{1 - q} .
\]
\end{example}

\begin{example}[Desimal yang periodik adalah deret geometri]\label{ex:b1:series:periodicdecimal}
Bilangan apakah $0.363636\dots$? Penulisannya sendiri sebuah \omterm{def:b1:series:def}{deret}:
\[
0.\overline{36} = \sum_{k=1}^{\infty} \frac{36}{100^k}
= 36\cdot\frac{1/100}{1 - 1/100} = \frac{36}{99} =
\frac{4}{11} ,
\]
menurut jumlah geometrinya dengan $q = \frac{1}{100}$. Secara umum sebuah
blok $B$ berisi $p$ angka yang berulang selamanya bernilai
$\frac{B}{10^p - 1}$ --- yaitu mekanisme di balik kriteria keperiodikan
pada \cref{pb:b1:reals:1}, yang akhirnya dinyatakan bahasa bab ini
dalam satu baris: bahwa \omterm{pb:b1:reals:1}{ekspansi desimal} merupakan \omterm{def:b1:series:def}{deret} yang konvergen,
yang periodik pada akhirnya tepat ketika jumlahnya rasional.
Jadi mesin angka Bab 10, yang dibangun di sana dengan \omterm{def:b1:reals:bounds}{supremum}
telanjang, sebenarnya teori \omterm{def:b1:series:def}{deret} yang menyamar.
\end{example}

\begin{proposition}[Fakta yang pertama]\label{prop:b1:series:first}
\begin{enumerate}
  \item Jika $\sum u_n$ konvergen, maka $u_n \to 0$. (Adapun konversnya
        \emph{salah}: yaitu \omterm{def:b1:series:def}{deret} harmoniknya.)
  \item Kelinearan: bahwa \omterm{def:b1:series:def}{deret} yang konvergen terjumlahkan dan terskalakan, dengan
        jumlah yang diharapkan.
  \item (Teleskopis\index{deret teleskopis}) $\sum (v_{n+1} -
        v_n)$ konvergen jika dan hanya jika $(v_n)$ konvergen, dengan jumlah $\lim v_n -
        v_0$.
  \item Mengubah suku yang berhingga banyak tak memengaruhi kekonvergenannya
        (melainkan hanya jumlahnya).
\end{enumerate}
\end{proposition}

\begin{proof}
(1) Karena $u_N = S_N - S_{N-1} \to S - S = 0$. Adapun \omterm{def:b1:series:def}{deret} harmoniknya mempunyai
$u_n = \frac1n \to 0$ namun divergen (\cref{exo:b1:seq:5}). (2)
Lewat operasi pada limitnya. (3) $S_N = v_{N+1} - v_0$. (4) Karena jumlah
parsialnya berubah sebesar besaran yang akhirnya konstan.
\end{proof}

\begin{example}[Merencanakan angka dengan sisa geometrinya]\label{ex:b1:series:georemainder}
Untuk $\abs q < 1$ sisa \omterm{ex:b1:series:geometric}{deret geometrinya} bersifat
eksplisit:
\[
R_N = \sum_{n = N+1}^{\infty} q^n = \frac{q^{N+1}}{1 - q} .
\]
Ini mengubah sasaran ketelitian menjadi cacah suku sebelum perhitungan
apa pun. Untuk menilai $\sum_{n\geq0} \bigl(\frac13\bigr)^n =
\frac32$ dalam jarak $10^{-10}$: kita memerlukan $\frac{(1/3)^{N+1}}{2/3} \leq
10^{-10}$, yakni $3^{N} \geq \frac{3}{2}\cdot 10^{10}$, yakni $N
\geq 22$ (karena $3^{22} \approx 3.1\cdot10^{10}$): jadi dua puluh tiga
suku, yang diketahui di muka. Adapun setiap taksiran berlaju geometri pada
soal akhir pekannya (yaitu \omterm{def:b1:series:def}{deret} $\frac13$ bagi $\ln 2$, dan
\omterm{def:b1:functions:arc}{arkus tangen} Machin pada \cref{pb:b1:taylor:1}) merupakan anggaran dua baris
ini dalam pakaian profesional.
\end{example}

\begin{example}[Teleskop yang lebih panjang]\label{ex:b1:series:telescope3}
Hitunglah $\sum_{n\geq1} \frac{1}{n(n+1)(n+2)}$. Lewat pecahan
parsialnya (\cref{ch:b1:fractions}):
\[
\frac{1}{n(n+1)(n+2)}
= \frac{1/2}{n} - \frac{1}{n+1} + \frac{1/2}{n+2}
= \frac12\Bigl(\frac{1}{n(n+1)} - \frac{1}{(n+1)(n+2)}\Bigr),
\]
yang di situ bentuk keduanya --- yaitu selisih nilai berurutan
$w_n = \frac{1}{n(n+1)}$ --- merupakan yang teleskopis. Sehingga
\[
\sum_{n=1}^{N} \frac{1}{n(n+1)(n+2)}
= \frac12\Bigl(w_1 - w_{N+1}\Bigr)
= \frac12\Bigl(\frac12 - \frac{1}{(N+1)(N+2)}\Bigr)
\longrightarrow \frac14 .
\]
Inti gagasan penutupnya: bahwa pecahan parsial bersuku tiga jarang
berteleskop sebagaimana ditulis; jadi kelompokkanlah ulang menjadi selisih $w_n -
w_{n+1}$ lebih dulu --- dan hadiahnya bukan hanya kekonvergenan melainkan
jumlahnya yang persis, yang tak pernah diberikan uji perbandingan mana pun.
\end{example}

\section{Deret bersuku taknegatif}

\begin{theorem}[Jumlah parsial yang terbatas]\label{thm:b1:series:positive}
Jika $u_n \geq 0$ untuk setiap $n$, maka jumlah parsialnya naik, sehingga: $\sum
u_n$ konvergen $\iff$ jumlah parsialnya terbatas di atas. Jadi berlakulah
\emph{uji perbandingan}: bahwa jika $0 \leq u_n \leq v_n$ untuk setiap (yang besar)
$n$,
\[
\sum v_n \text{ konvergen} \implies \sum u_n \text{ konvergen},
\qquad
\sum u_n \text{ divergen} \implies \sum v_n \text{ divergen}.
\]
Dan \emph{uji kesetaraan}: bahwa jika $u_n \sim v_n$ dengan $v_n \geq
0$, maka kedua \omterm{def:b1:series:def}{deretnya} bersifat sama.
\end{theorem}

\begin{proof}
Lewat teorema limit monoton (\cref{thm:b1:seq:monotone}) untuk butir
pertamanya; dan lewat perbandingan jumlah parsialnya untuk yang kedua. Adapun kesetaraannya: untuk
$n$ yang besar, $\frac12 v_n \leq u_n \leq 2 v_n$ (menurut definisi $\sim$
dengan $\varepsilon = \frac12$), lalu perbandingannya berlaku ke kedua arahnya.
\end{proof}

\begin{example}[Sebuah kesetaraan yang membuktikan kedivergenan]\label{ex:b1:series:equivdiverge}
Bagaimanakah sifat $\sum_{n\geq1} n\sin\dfrac{1}{n^2}$? Karena
$\frac{1}{n^2} \to 0$ dan $\sin h \sim h$ di $0$:
\[
n\sin\frac{1}{n^2} \;\sim\; n\cdot\frac{1}{n^2} = \frac1n ,
\]
lalu uji kesetaraannya memindahkan kedivergenan \omterm{def:b1:series:def}{deret}
harmoniknya: jadi divergen --- meskipun sukunya menuju
$0$. Jadi satu ekspansi, satu skala, satu vonis; dan pola dua langkah
yang sama (yaitu kesetaraan, lalu pencarian Riemann atau geometri) memutuskan
keempat \omterm{def:b1:series:def}{deret} pada \cref{exo:b1:series:3}.
\end{example}

\begin{example}[Uji kesetaraan dalam satu baris]\label{ex:b1:series:equivtest}
Bagaimanakah sifat $\sum_{n \geq 1} \frac{\sqrt{n+1} - \sqrt n}{n}$?
Sekawankanlah pembilangnya:
\[
\frac{\sqrt{n+1} - \sqrt n}{n}
= \frac{1}{n\,(\sqrt{n+1} + \sqrt n)}
\sim \frac{1}{2\,n^{3/2}} ,
\]
yaitu skala Riemann yang konvergen (dengan $\alpha = \frac32 > 1$): sehingga \omterm{def:b1:series:def}{deretnya}
konvergen. Jadi seluruh keputusannya memerlukan satu kesetaraan dan satu pencarian
--- asalkan sukunya taknegatif, dan memang demikian. Inti
gagasan penutupnya: bahwa bagi \omterm{def:b1:series:def}{deret} yang positif, seluruh teori
kekonvergenannya merupakan \emph{kamus skala} ($n^{-\alpha}$, $q^n$,
$\frac{1}{n(\ln n)^\alpha}$) ditambah izin mengganti sebuah suku
dengan kesetaraannya; sedangkan kerja analitiknya ada pada asimtotiknya
(\cref{ch:b1:taylor}), tak pernah pada penjumlahannya.
\end{example}

\begin{theorem}[Perbandingan integral; deret Riemann]\label{thm:b1:series:riemann}
Misalkan $f$ \omterm{def:b1:continuity:continuous}{kontinu}, taknegatif dan \emph{menurun} pada
$\intco{1}{+\infty}$. Maka
\[
\int_1^{N+1} f(t)\,\dd t \;\leq\; \sum_{n=1}^{N} f(n) \;\leq\; f(1)
+ \int_1^{N} f(t)\,\dd t ,
\]
sehingga $\sum f(n)$ konvergen jika dan hanya jika $\bigl(\int_1^x f\bigr)$ terbatas.
Khususnya, untuk $\alpha \in \R$:\index{deret Riemann}
\[
\sum_{n \geq 1} \frac{1}{n^\alpha} \text{ konvergen}
\iff \alpha > 1,
\]
dan $\sum_{n=1}^{N} \frac1n = \ln N + O(1)$.
\end{theorem}

\begin{proof}
Untuk $n \leq t \leq n+1$, kemonotonannya memberikan $f(n+1) \leq f(t)
\leq f(n)$; lalu mengintegralkannya atas $\intcc{n}{n+1}$ (yaitu ruas
berpanjang $1$):
\[
f(n+1) \;\leq\; \int_n^{n+1} f(t)\,\dd t \;\leq\; f(n) .
\]
Lalu menjumlahkan ketaksamaan kanannya untuk $n = 1, \dots, N-1$ memberikan
$\int_1^{N} f \leq \sum_{n=1}^{N-1} f(n)$, sehingga apitan atasnya diperoleh
setelah menambahkan $f(N) \leq f(1)$; sedangkan menjumlahkan yang kiri untuk
$n = 1, \dots, N$ memberikan $\sum_{n=2}^{N+1} f(n) \leq
\int_1^{N+1} f$, yang setelah diindeks ulang menjadi apitan bawahnya.
Adapun kekonvergenannya: jumlah parsialnya dan \omterm{thm:b1:integration:def}{integral} $\int_1^x f$
saling membatasi dalam jarak konstanta $f(1)$, dan keduanya
tidak turun, sehingga yang satu terbatas jika dan hanya jika yang lain demikian
(\cref{thm:b1:series:positive}). Untuk $f(t)
= t^{-\alpha}$ (dengan $\alpha \neq 1$): $\int_1^x t^{-\alpha}\dd t =
\frac{x^{1-\alpha} - 1}{1 - \alpha}$, yang terbatas jika dan hanya jika $\alpha > 1$; sedangkan untuk
$\alpha = 1$ \omterm{thm:b1:integration:def}{integralnya} $\ln x \to \infty$, dan apitannya
memberikan $\ln(N+1) \leq H_N \leq 1 + \ln N$. Adapun untuk $\alpha \leq 0$
sukunya tak menuju $0$.
\end{proof}

\begin{example}[Tumpukan harmoniknya]\label{ex:b1:series:harmonicstack}
Berapa banyak suku yang harus dikumpulkan \omterm{def:b1:series:def}{deret} harmoniknya untuk melewati $20$?
Apitan $\ln(N+1) \leq H_N \leq 1 + \ln N$ menjawabnya tanpa
penjumlahan sama sekali: karena $H_N \geq 20$ menuntut $1 + \ln N \geq 20$,
yakni $N \geq \eu^{19} \approx 1.8\cdot10^{8}$, dan terjamin
begitu $\ln(N + 1) \geq 20$, yakni $N \approx \eu^{20} \approx
4.9\cdot10^{8}$. (Adapun soal akhir pekannya mempertajam ini menjadi $N \approx
\eu^{20 - \gamma} \approx 2.7\cdot10^{8}$ lewat \omterm{pb:b1:series:1}{konstanta Euler}.)
Inti gagasan penutupnya: bahwa perbandingan \omterm{thm:b1:integration:def}{integralnya} tak sekadar
memutuskan kekonvergenan --- ia \emph{melokalkan} jumlah parsialnya dengan
ketelitian logaritmis, sehingga mengubah perhitungan yang tanpa harapan
(yaitu ratusan juta suku) menjadi taksiran dua baris.
\end{example}

\begin{theorem}[Uji nisbah (d'Alembert)]\label{thm:b1:series:ratio}
Misalkan $u_n > 0$ dengan $\frac{u_{n+1}}{u_n} \to \ell$.\index{uji nisbah}
\begin{itemize}
  \item Jika $\ell < 1$: maka $\sum u_n$ konvergen;
  \item jika $\ell > 1$: maka $u_n \to +\infty$, jadi divergen;
  \item jika $\ell = 1$: maka tak ada kesimpulannya (karena $\sum \frac1n$ divergen, sedangkan $\sum
        \frac{1}{n^2}$ konvergen).
\end{itemize}
\end{theorem}

\begin{proof}
Jika $\ell < 1$, tetapkanlah $q \in \intoo{\ell}{1}$: maka di luar suatu $N$,
$u_{n+1} \leq q\,u_n$, sehingga $u_n \leq u_N q^{\,n-N}$ lewat induksi:
jadi perbandingan dengan \omterm{ex:b1:series:geometric}{deret geometri}. Jika $\ell > 1$: maka di luar suatu $N$
barisan $(u_n)$ naik, sehingga ia tak dapat menuju $0$
(karena limitnya, bila ada, bernilai $\geq u_N > 0$); jadi menurut
\cref{prop:b1:series:first}~(1), divergen --- dan sesungguhnya $u_n
\geq u_N q^{n-N}$ dengan $q > 1$ memberikan $u_n \to \infty$.
\end{proof}

\begin{example}\label{ex:b1:series:ratio}
$\sum \frac{x^n}{n!}$ konvergen untuk setiap $x > 0$: karena nisbahnya
$\frac{x}{n+1} \to 0$. Adapun jumlahnya $\eu^x$: karena menurut Taylor--Lagrange
(\cref{thm:b1:taylor:lagrange}) pada $\intcc{0}{x}$,
\[
\Bigl| \eu^x - \sum_{k=0}^{n} \frac{x^k}{k!} \Bigr|
\leq \eu^{x}\, \frac{x^{n+1}}{(n+1)!} \xrightarrow[n\to\infty]{} 0 ,
\]
dengan batasnya menuju $0$ karena faktorialnya mendominasi
(\cref{exo:b1:integration:9}~(1) memakai fakta yang sama). Adapun argumen yang
sama menjumlahkan \omterm{def:b1:series:def}{deret} $\sin$, $\cos$, $\sinh$, dan $\cosh$ pada seluruh
$\R$.
\end{example}

\begin{example}[Uji nisbahnya mencukupi, tetapi tak diperlukan]\label{ex:b1:series:rationotnecessary}
Misalkan $u_n = 2^{-n}$ untuk $n$ yang genap dan $u_n = 2^{-n-2}$ untuk $n$ yang
ganjil. Maka nisbah berurutannya berayun di antara $\frac{1}{8}$ dan
$\frac12\cdot4 = 2$, sehingga $\frac{u_{n+1}}{u_n}$ tak berlimit dan
d'Alembert bungkam --- padahal $u_n \leq 2^{-n}$ dan uji
perbandingannya menuntaskan kekonvergenannya seketika. Jadi hipotesis ujinya (yaitu bahwa
nisbahnya \emph{konvergen}) merupakan pembatasan yang sungguhan: ia cocok bagi suku
dengan satu struktur perkalian yang dominan (yaitu faktorial dan pangkat),
dan gagal pada apa pun yang bernapas. Jadi ketika nisbahnya berkelakuan buruk,
mundurlah ke perbandingan terhadap selubung geometri --- dan hanya itulah
\omterm{thm:b1:series:ratio}{uji nisbahnya} sejak semula, sebagaimana diperlihatkan buktinya.
\end{example}

\begin{example}[Uji nisbah pada pertempuran faktorial]\label{ex:b1:series:ratiobinom}
Bagaimanakah sifat $\sum_{n\geq0} \dfrac{(n!)^2}{(2n)!}$ (yaitu kebalikan
\omterm{def:b1:counting:objects}{koefisien binomial} pusatnya, sampai faktor $n + 1$)?
Nisbahnya meruntuhkan faktorialnya:
\[
\frac{u_{n+1}}{u_n}
= \frac{((n+1)!)^2}{(n!)^2}\cdot\frac{(2n)!}{(2n+2)!}
= \frac{(n+1)^2}{(2n+1)(2n+2)}
\longrightarrow \frac14 < 1 :
\]
jadi konvergen, dengan kelonggaran --- karena sukunya meluruh pada dasarnya
seperti $4^{-n}$, yang selaras dengan $\binom{2n}{n} \geq
\frac{4^n}{2n+1}$ dari \cref{pb:b1:integration:1}. Inti gagasan
penutupnya: bahwa hasil bagi faktorial persis merupakan apa yang dicerna \omterm{thm:b1:series:ratio}{uji nisbahnya}
--- karena setiap faktorialnya saling meniadakan menjadi fungsi rasional dalam
$n$, yang limitnya terbaca dari suku utamanya.
\end{example}

\section{Kekonvergenan mutlak; deret berselang-seling}

\begin{theorem}[Kekonvergenan mutlak]\label{thm:b1:series:absolute}
Jika $\sum \abs{u_n}$ konvergen (yaitu \emph{kekonvergenan
mutlak}\index{kekonvergenan mutlak}), maka $\sum u_n$ konvergen,
dan $\bigl|\sum u_n\bigr| \leq \sum \abs{u_n}$. Ini berlaku bagi suku yang real
maupun kompleks.
\end{theorem}

\begin{proof}
Jumlah parsialnya memenuhi, untuk $M > N$ (menurut kriteria Cauchy,
\cref{thm:b1:seq:complete}):
\[
\abs{S_M - S_N} = \Bigl| \sum_{n=N+1}^{M} u_n \Bigr|
\leq \sum_{n=N+1}^{M} \abs{u_n},
\]
yang bernilai kecil untuk $N$ yang besar karena jumlah parsial $\sum\abs{u_n}$
membentuk \omterm{def:b1:seq:cauchy}{barisan Cauchy}. Jadi $(S_N)$ bersifat Cauchy, sehingga konvergen. Adapun
ketaksamaannya beralih ke limitnya dari ketaksamaan segitiga yang hingga.
\end{proof}

\begin{example}[Kekonvergenan mutlak, yang real dan yang kompleks]\label{ex:b1:series:absexamples}
Untuk $\sum_{n\geq1} \frac{\sin n}{n^2}$: sukunya berubah tanda
secara tak beraturan (memang $(\sin n)$ \omterm{def:b1:topology:dense}{padat} di $\intcc{-1}{1}$,
\cref{exo:b1:seq:12}), dan tak ada struktur berselang-seling yang
terlihat. Namun \omterm{thm:b1:series:absolute}{kekonvergenan mutlaknya} menyelamatkan segalanya sekaligus:
karena $\bigl|\frac{\sin n}{n^2}\bigr| \leq \frac{1}{n^2}$, yaitu
skala yang konvergen, sehingga \omterm{def:b1:series:def}{deretnya} konvergen. Adapun perisai yang sama bekerja
di atas $\C$: $\sum_{n\geq1}\frac{\eu^{\iu n}}{n^2}$ konvergen
karena $\bigl|\frac{\eu^{\iu n}}{n^2}\bigr| = \frac{1}{n^2}$
--- jadi pola tandanya, bahkan yang berdimensi dua, tak relevan
begitu modulusnya terjumlahkan. Inti gagasan penutupnya: bahwa \omterm{thm:b1:series:absolute}{kekonvergenan
mutlak} merupakan satu-satunya perkakas bab ini yang tak pernah menanyakan bagaimana
tandanya tertata; jadi cobalah ia lebih dulu
(\cref{met:b1:series:decide}), lalu sisakan uji yang halus
bagi \omterm{def:b1:series:def}{deret} yang menggagalkannya.
\end{example}

\begin{theorem}[Uji deret berselang-seling]\label{thm:b1:series:alternating}
Misalkan $(a_n)$ menurun dengan $a_n \to 0$. Maka \omterm{thm:b1:series:alternating}{deret
berselang-seling} $\sum (-1)^n a_n$ konvergen; jumlahnya terletak di antara sebarang dua
jumlah parsial yang berurutan, dan\index{deret berselang-seling}
\[
\abs{R_N} = \Bigl| \sum_{n > N} (-1)^n a_n \Bigr| \leq a_{N+1} .
\]
\end{theorem}

\begin{proof}
Jumlah parsial genap dan ganjilnya berdampingan: karena $S_{2p+2} - S_{2p} =
a_{2p+2} - a_{2p+1} \leq 0$ (menurut penurunannya), $S_{2p+1} - S_{2p-1} =
a_{2p} - a_{2p+1} \geq 0$ (menurut kenaikannya), dan $S_{2p} - S_{2p+1} =
a_{2p+1} \to 0$. Jadi menurut \cref{thm:b1:seq:adjacent} keduanya berbagi limit
$S$, yang oleh kriteria dua \omterm{def:b1:seq:subsequence}{subbarisan}
(\cref{prop:b1:seq:subsequences}) menjadi limit $(S_N)$;
lebih lanjut $S$ terperangkap di antara jumlah parsial yang berurutan, dan
$\abs{S - S_N}$ paling banyak sebesar celah ke yang berikutnya, yaitu $a_{N+1}$.
\end{proof}

\begin{example}[Deret harmonik berselang-seling]\label{ex:b1:series:ln2}
$\sum_{n \geq 1} \frac{(-1)^{n-1}}{n}$ konvergen (menurut uji berselang-selingnya)
tetapi tak mutlak (karena \omterm{def:b1:series:def}{deret} harmoniknya). Adapun jumlahnya $\ln 2$: karena dari
kesamaan geometri yang hingga $\frac{1}{1+t} = \sum_{k=0}^{n-1} (-t)^k +
\frac{(-t)^n}{1+t}$, integralkanlah atas $\intcc{0}{1}$:
\[
\ln 2 = \sum_{k=1}^{n} \frac{(-1)^{k-1}}{k}
+ (-1)^n \int_0^1 \frac{t^n}{1+t}\,\dd t ,
\qquad
0 \leq \int_0^1 \frac{t^n}{1+t}\,\dd t \leq \frac{1}{n+1} \to 0 .
\]
Kekonvergenannya lamban dengan menyakitkan (karena $R_N \approx \frac{1}{N}$) ---
sebab \omterm{thm:b1:series:alternating}{deret berselang-seling} konvergen lewat peniadaan, bukan lewat kekecilan.
\end{example}

\begin{omfigure}
\begin{tikzpicture}[x=0.78cm, y=7.2cm]
  \draw[-Stealth] (0.3,0.42) -- (11,0.42)
    node[below, font=\small] {$N$};
  \draw (1,0.415) -- (1,0.425) node[below=1pt, font=\small]
    {$1$};
  \draw (5,0.415) -- (5,0.425) node[below=1pt, font=\small]
    {$5$};
  \draw (10,0.415) -- (10,0.425) node[below=1pt, font=\small]
    {$10$};
  \draw[gray, dashed] (0.3,0.6931) -- (10.7,0.6931)
    node[right, font=\small, text=black] {$\ln 2$};
  \node[left, font=\small] at (0.3,1) {$1$};
  \node[left, font=\small] at (0.3,0.5) {$0.5$};
  \draw[omDef, thick]
    plot coordinates {(1,1) (2,0.5) (3,0.8333) (4,0.5833)
    (5,0.7833) (6,0.6167) (7,0.7595) (8,0.6345)
    (9,0.7456) (10,0.6456)};
  \fill[omDef] (1,1) circle (1.5pt);
  \fill[omDef] (2,0.5) circle (1.5pt);
  \fill[omDef] (3,0.8333) circle (1.5pt);
  \fill[omDef] (4,0.5833) circle (1.5pt);
  \fill[omDef] (5,0.7833) circle (1.5pt);
  \fill[omDef] (6,0.6167) circle (1.5pt);
  \fill[omDef] (7,0.7595) circle (1.5pt);
  \fill[omDef] (8,0.6345) circle (1.5pt);
  \fill[omDef] (9,0.7456) circle (1.5pt);
  \fill[omDef] (10,0.6456) circle (1.5pt);
\end{tikzpicture}

{\small Jumlah parsial $S_N$ \omterm{def:b1:series:def}{deret} harmonik berselang-seling
$1 - \frac12 + \frac13 - \cdots$ melompati limitnya
$\ln 2$ pada setiap langkahnya: yaitu jumlah ganjilnya dari atas, dan jumlah genapnya dari
bawah, dengan setiap lompatannya berukuran $\frac{1}{N+1}$. Jadi apitannya merupakan
bukti \cref{thm:b1:series:alternating} yang dibuat kasatmata --- dan
lambannya penutupan penjepitnya ($\abs{S_N - \ln 2} \approx
\frac{1}{2N}$, menurut soal akhir pekan \cref{pb:b1:series:1}) adalah sebabnya
tak seorang pun menghitung $\ln 2$ dengan cara ini.}
\end{omfigure}

\begin{remark}[Jebakan yang lazim dengan deret]
(i) \emph{Uji kesetaraannya memerlukan tanda}: misalkan $v_n =
\frac{(-1)^n}{\sqrt n}$ dan $u_n = v_n + \frac1n$. Maka
$\frac{u_n}{v_n} = 1 + \frac{(-1)^n}{\sqrt n} \to 1$, sehingga $u_n
\sim v_n$; namun $\sum v_n$ konvergen (menurut uji berselang-selingnya) sedangkan
$\sum u_n = \sum v_n + \sum \frac1n$ divergen. Jadi kesetaraan
mengendalikan \emph{ukuran} sukunya, dan bagi \omterm{def:b1:series:def}{deret} yang bertanda ukurannya
bukan takdir --- sehingga ujinya dinyatakan, dan benar, hanya bagi suku yang
(akhirnya) taknegatif. (ii) \emph{$u_n \to 0$
tak membuktikan apa pun}: karena \omterm{def:b1:series:def}{deret} harmoniknya merupakan contoh penyangkal
yang abadi; dan arah konversnya
(\cref{prop:b1:series:first}~(1)) hanyalah uji kedivergenan
yang cepat. (iii) \emph{Limit nisbah $1$ itu kebungkaman, bukan
kekonvergenan}: karena baik $\sum\frac1n$ maupun $\sum\frac{1}{n^2}$ bernisbah
$\to 1$; jadi berpindahlah ke skala Riemann atau perbandingan \omterm{thm:b1:integration:def}{integral}.
(iv) \emph{Berselang-seling memerlukan penurunan}: karena $\sum
\frac{(-1)^n}{n + (-1)^n}$ tampak berselang-seling padahal ditangani
hanya lewat ekspansi (\cref{exo:b1:series:5}); dan soal akhir pekan
\cref{ch:b1:taylor} (pertanyaan 23 di sana) menunjukkan bahwa ujinya dapat
gagal sama sekali tanpa kemonotonannya. (v) \emph{Pengelompokan dan
penyusunan ulang tak cuma-cuma}: karena menyisipkan tanda kurung tak berbahaya bagi
\omterm{def:b1:series:def}{deret} yang konvergen tetapi dapat menciptakan kekonvergenan dari kedivergenan
(yaitu $1 - 1 + 1 - \cdots$ yang dikelompokkan berpasangan), dan penyusunan ulang dapat
mengubah jumlahnya sendiri --- yaitu drama yang dipentaskan pada
soal akhir pekan bab ini (\cref{pb:b1:series:1}).
\end{remark}

\begin{method}[Memutuskan sifat sebuah deret]\label{met:b1:series:decide}
\begin{enumerate}
  \item Apakah $u_n \to 0$? Jika tidak, divergen, berhenti.
  \item Untuk suku yang taknegatif: carilah kesetaraan bagi $u_n$ (lewat ekspansi,
        \cref{ch:b1:taylor}!), lalu bandingkanlah dengan skala Riemann atau
        geometri; adapun faktorial dan pangkat memanggil \omterm{thm:b1:series:ratio}{uji nisbahnya};
        sedangkan $f(n)$ yang menurun memanggil perbandingan \omterm{thm:b1:integration:def}{integralnya}.
  \item Bila tandanya berubah-ubah: cobalah \omterm{thm:b1:series:absolute}{kekonvergenan mutlaknya} lebih dulu; dan jika gagal,
        uji berselang-selingnya (periksalah \emph{penurunannya} dengan saksama); adapun di luar
        itu, perkakas tahun kedua.
\end{enumerate}
\end{method}

\begin{example}[Penyebut yang ganjil, separuh teleskopnya]\label{ex:b1:series:oddtelescope}
Hitunglah $\sum_{n \geq 1} \dfrac{1}{4n^2 - 1}$. Lewat pecahan
parsialnya: $\frac{1}{(2n-1)(2n+1)} =
\frac12\bigl(\frac{1}{2n-1} - \frac{1}{2n+1}\bigr)$, sehingga
\[
\sum_{n=1}^{N} \frac{1}{4n^2 - 1}
= \frac12\Bigl(1 - \frac{1}{2N+1}\Bigr)
\longrightarrow \frac12 .
\]
Bandingkanlah dengan $\sum \frac{1}{n(n+1)} = 1$
(\cref{exo:b1:series:1}): yaitu kerangka teleskopis yang sama, tetapi
suku berurutannya di sini berjarak dua pada bilangan ganjilnya, dan
faktor $\frac12$-nya mencatat langkahnya. Inti gagasan penutupnya:
bahwa teleskopis merupakan perubahan sudut pandang, bukan muslihat --- karena setiap kali
suku umumnya berupa selisih $w_n - w_{n+1}$ atas sebuah barisan
yang berlimit, jumlahnya adalah $w_1 - \lim w$, yang persis
\cref{prop:b1:series:first}~(3).
\end{example}

\begin{remark}[Pipa analisisnya, bila ditengok ke belakang]
Bab inilah tempat analisis jilid ini bertemu, dan setiap
ujinya menyebut leluhurnya. Jumlah parsial yang terbatas merupakan teorema limit
monoton (\cref{ch:b1:seq}), yang sendirinya aksioma kelengkapan
pada \cref{ch:b1:reals}; \omterm{thm:b1:series:absolute}{kekonvergenan mutlaknya} merupakan kriteria
Cauchy; uji \omterm{thm:b1:integration:def}{integralnya} merupakan apitan luas pada
\cref{ch:b1:integration}; kesetaraan suku umumnya merupakan
ekspansi pada \cref{ch:b1:taylor}; dan teorema berselang-selingnya merupakan
lema \omterm{thm:b1:seq:adjacent}{barisan berdampingan} dalam pakaian Minggunya. Jadi bila dibaca
mundur, pipanya menjelaskan setiap babnya \emph{untuk} apa
ada --- dan soal akhir pekan yang dirangkaikan di sepanjangnya
(yaitu angka $b$-adik, Cesàro--Stolz, mesin keirasionalan, dan
\omterm{pb:b1:series:1}{konstanta Euler}) merupakan beberapa gagasan yang sama yang bertemu pada ketinggian
yang kian tinggi. Adapun aljabar linear yang menyusul mengubah
pokoknya, bukan standarnya: karena kebiasaan \omterm{def:b1:logic:statement}{pernyataan} yang persis dengan
galat yang bersertifikat bertahan melewati perpindahan dari limit ke dimensi.
\end{remark}

\begin{remark}[Ke mana deretnya berlanjut]
Bab ini menutup analisis jilid ini dan membuka tiga
pintu. Pada jilid Tahun ke-2, \omterm{def:b1:series:def}{deretnya} memperoleh sebuah variabel ($\sum
a_n x^n$: yaitu \omterm{def:b1:series:def}{deret} pangkat, dengan jari-jari kekonvergenannya) lalu
kemudian sebuah teori bernilai fungsi (yaitu \omterm{def:b1:series:def}{deret} Fourier); dan dikotomi
mutlak lawan bersyaratnya, yang didramatiskan pada
soal akhir pekan di bawah, menjadi batu penjuru keduanya. Pada
peluang (jilid Tahun ke-3), ekspektasi variabel acak yang
diskret \emph{adalah} \omterm{def:b1:series:def}{deret}, dan kekonvergenan mutlaknyalah yang
membuatnya terdefinisi dengan baik. Sedangkan \omterm{thm:b1:series:riemann}{deret Riemann} $\sum n^{-s}$,
bila didorong ke $s$ yang kompleks, menjadi fungsi zeta --- yaitu satu
\omterm{def:b1:series:def}{deret} yang paling banyak ditelaah dalam matematika.
\end{remark}

\section{Latihan}

\begin{exercise}[$\star$]\label{exo:b1:series:1}
Sifat (dan jumlahnya, bila teleskopis) dari:
\[
\sum_{n\geq1} \frac{1}{n(n+1)},
\qquad
\sum_{n\geq2} \ln\Bigl(1 - \frac{1}{n^2}\Bigr),
\qquad
\sum_{n\geq0} \frac{3^n + 4^n}{5^n} .
\]
\end{exercise}

\begin{exercise}[$\star$]\label{exo:b1:series:2}
Sifat dari: $\;\sum \dfrac{n^2}{2^n}$; $\;\sum \dfrac{n!}{n^n}$;
$\;\sum \dfrac{2^n\,n!}{n^n}$; $\;\sum \dfrac{3^n\,n!}{n^n}$.
\emph{(Lewat \omterm{thm:b1:series:ratio}{uji nisbahnya}; ingatlah $\bigl(1 + \frac1n\bigr)^n \to \eu$.)}
\end{exercise}

\begin{exercise}[$\star$]\label{exo:b1:series:3}
Sifat dari: $\;\sum \sin\dfrac{1}{n^2}$; $\;\sum
\Bigl(1 - \cos\dfrac1n\Bigr)$; $\;\sum \dfrac{1}{\sqrt{n(n+1)}}$;
$\;\sum \dfrac{\ln n}{n^2}$ \emph{(bandingkanlah dengan $n^{-3/2}$)}.
\end{exercise}

\begin{exercise}[$\star$]\label{exo:b1:series:4}
Buktikan bahwa $\sum_{n\geq1} \frac{1}{n^2}$ konvergen dengan jumlah $\leq 2$,
dengan memakai $\frac{1}{n^2} \leq \frac{1}{n(n-1)}$ untuk $n \geq 2$ beserta batas
teleskopisnya.
\end{exercise}

\begin{exercise}[$\star\star$]\label{exo:b1:series:5}
Sifat dari $\;\sum \dfrac{(-1)^n}{\sqrt n}$, dari $\;\sum
\dfrac{(-1)^n}{n + (-1)^n}$ \emph{(ekspansikanlah: karena uji berselang-selingnya tak
berlaku secara langsung --- mengapa?)}, dan dari $\;\sum
\sin\bigl(\pi\sqrt{n^2+1}\,\bigr)$ \emph{(reduksikanlah modulo $\pi$:
$\sqrt{n^2+1} = n + \frac{1}{2n} + O(n^{-3})$)}.
\end{exercise}

\begin{exercise}[$\star\star$]\label{exo:b1:series:6}
Untuk $\alpha > 0$ yang mana $\sum_{n \geq 2}
\dfrac{1}{n (\ln n)^{\alpha}}$ konvergen? \emph{(Lewat perbandingan \omterm{thm:b1:integration:def}{integral};
sulihkanlah $u = \ln t$.)}
\end{exercise}

\begin{exercise}[$\star\star$]\label{exo:b1:series:7}
Misalkan $u_n = \dfrac{1}{n} - \ln\Bigl(1 + \dfrac1n\Bigr)$. Buktikan bahwa
$0 \leq u_n \leq \dfrac{1}{2n^2}$, bahwa $\sum u_n$ konvergen, lalu
simpulkanlah keberadaan \omterm{pb:b1:series:1}{konstanta Euler}\index{konstanta Euler}:
\[
\gamma = \lim_{N \to \infty} \Bigl( \sum_{n=1}^{N} \frac 1n - \ln N
\Bigr) .
\]
\end{exercise}

\begin{exercise}[$\star\star$]\label{exo:b1:series:8}
Hitunglah jumlahnya
\[
\sum_{n=1}^{\infty} \frac{1}{n(n+2)}
\qquad\text{dan}\qquad
\sum_{n=0}^{\infty} \frac{n}{2^n} .
\]
\emph{(Untuk yang pertama: lewat pecahan parsialnya. Untuk yang kedua: hitunglah
$\sum_{n=1}^{N} n x^{n-1}$ dalam bentuk \omterm{def:b1:topology:closed}{tertutup} lalu biarkan $N \to \infty$ di
$x = \frac12$.)}
\end{exercise}

\begin{exercise}[$\star\star\star$]\label{exo:b1:series:9}
(Pemampatan Cauchy) Misalkan $(u_n)$ taknegatif dan menurun.
Buktikan bahwa
\[
\sum_{n \geq 1} u_n \text{ konvergen}
\iff
\sum_{k \geq 0} 2^k\, u_{2^k} \text{ konvergen},
\]
dengan membandingkan paket suku di antara pangkat $2$ yang berurutan.
Lalu pulihkanlah darinya kriteria Riemann beserta
\cref{exo:b1:series:6}.
\end{exercise}

\begin{exercise}[$\star\star\star$]\label{exo:b1:series:10}
Dengan memakai kesamaan \omterm{thm:b1:integration:def}{integral} pada \cref{ex:b1:series:ln2} yang disesuaikan bagi
$\frac{1}{1+t^2}$, buktikanlah rumus Leibniz
\[
\frac{\pi}{4} = \sum_{n=0}^{\infty} \frac{(-1)^n}{2n+1}
= 1 - \frac13 + \frac15 - \frac17 + \cdots
\]
dengan batas galatnya $\abs{R_N} \leq \frac{1}{2N+3}$.
\end{exercise}

\begin{exercise}[$\star\star$]\label{exo:b1:series:11}
Sifat dari $\displaystyle\sum_{n \geq 1} \frac{1}{n^{1 + 1/n}}$.
\emph{(Hitunglah limit $n^{1/n}$ lalu carilah kesetaraan bagi
suku umumnya: karena uji Riemann memerlukan eksponen yang \emph{tetap}.)}
\end{exercise}

\begin{exercise}[$\star\star\star$]\label{exo:b1:series:12}
Misalkan $(u_n)$ taknegatif dan \emph{menurun} dengan $\sum u_n$
yang konvergen. Buktikan bahwa $n\,u_n \to 0$ \emph{(batasilah $n\,u_{2n}$
oleh sebuah iris $\sum_{k=n+1}^{2n} u_k$ lalu pakailah kriteria
Cauchy)}. Tunjukkanlah bahwa konversnya gagal, dan bahwa
hipotesis kemonotonannya tak dapat dilepaskan.
\end{exercise}

\section{Soal: Konstanta Euler dan deret yang mengubah
jumlahnya}

\begin{problem}[{Soal akhir pekan --- $H_n = \ln n + \gamma +
\frac{1}{2n} + O(n^{-2})$, dan menyusun ulang $1 - \frac12 +
\frac13 - \dots$ menjadi $\frac{\ln 2}{2}$}]
\label{pb:b1:series:1}
Dua kisah berbagi \omterm{def:b1:series:def}{deret} harmoniknya. Pertama, pembukuan
yang persis atas kedivergenannya: bahwa $H_n - \ln n$ konvergen ke \omterm{pb:b1:series:1}{konstanta
Euler} $\gamma$ (\cref{exo:b1:series:7}), dan soal ini
mempertajam \omterm{def:b1:logic:statement}{pernyataannya} menjadi hukum dua sisi $\frac{1}{2(n+1)}
\leq H_n - \ln n - \gamma \leq \frac{1}{2n}$, yang mensertifikasi $\gamma
= 0.5772\dots$ dengan tangan. Kedua, skandal kekonvergenan
bersyarat: bahwa \omterm{def:b1:series:def}{deret} harmonik berselang-seling berjumlah $\ln 2$
(\cref{ex:b1:series:ln2}), namun \emph{suku yang sama, dalam
urutan yang berbeda}, berjumlah $\frac{\ln 2}{2}$ --- atau $\ln 2 +
\frac12\ln\frac pq$ untuk sebarang $p, q$, atau bahkan sebarang bilangan real
(menurut Riemann). Dan kedua kisahnya sesungguhnya satu: karena jumlah yang tersusun ulang itu
dihitung \emph{dengan} hukum $\gamma$-nya.\index{konstanta
Euler}\index{penyusunan ulang deret}

\medskip
\noindent\textbf{Bagian I --- $\gamma$, yang terapit.} Tetapkanlah $a_n = H_n
- \ln n$ dan $b_n = H_n - \ln(n+1)$.
\begin{enumerate}
  \item Dengan memakai $\frac{t}{1+t} \leq \ln(1 + t) \leq t$, tunjukkanlah bahwa
        $(a_n)$ menurun, $(b_n)$ menaik, dan bahwa keduanya
        berdampingan; sehingga limit bersamanya adalah $\gamma$, dengan $b_n \leq
        \gamma \leq a_n$ untuk setiap $n$.
  \item Tembakan numerik pertamanya: dari $H_{10} = 2.928968\dots$,
        apitlah $\gamma$ di antara $b_{10} = 0.5311$ dan $a_{10} =
        0.6264$. Seberapa besarkah $n$ yang diperlukan apitan yang kasar ini
        bagi empat desimal?
  \item Tunjukkanlah penyajian ekornya yang persis $a_n - \gamma =
        \sum_{k \geq n} w_k$ (yaitu limit jumlah parsialnya), dengan
        \[
        w_k = a_k - a_{k+1}
        = \ln\Bigl(1 + \frac1k\Bigr) - \frac{1}{k+1}
        = \int_k^{k+1} \frac{(k + 1 - t)}{t\,(k+1)}\,\dd t ,
        \]
        lalu simpulkanlah dari bentuk \omterm{thm:b1:integration:def}{integralnya} batas dua sisinya
        $\dfrac{1}{2(k+1)^2} \leq w_k \leq \dfrac{1}{2k(k+1)}$.
\end{enumerate}

\medskip
\noindent\textbf{Bagian II --- Hukum $\frac{1}{2n}$-nya.}
\begin{enumerate}[resume]
  \item Jumlahkanlah batas pertanyaan 3 (karena kedua sisinya berteleskop atau
        terbandingkan dengan teleskop) lalu simpulkanlah hukumnya:
        \[
        \frac{1}{2(n+1)} \;\leq\; H_n - \ln n - \gamma \;\leq\;
        \frac{1}{2n} \qquad (n \geq 1).
        \]
  \item Simpulkanlah $H_n = \ln n + \gamma + \frac{1}{2n} +
        O\bigl(\frac{1}{n^2}\bigr)$; lebih tepatnya, tunjukkanlah bahwa
        $\gamma_n = H_n - \ln n - \frac{1}{2n}$ memenuhi
        $-\frac{1}{2n(n+1)} \leq \gamma_n - \gamma \leq 0$.
  \item Sertifikasikanlah empat desimal dengan $n = 100$: diberikan $H_{100} =
        5.1873775\dots$, hitunglah $\gamma_{100} = 0.577207\dots$
        lalu simpulkanlah $\gamma = 0.5772 \pm 5\cdot10^{-5}$ (dengan nilai
        sejatinya $0.5772156\dots$).
  \item Dua dividen hukumnya, yang keduanya diperlukan nanti: bahwa ketika $m \to
        \infty$,
        \[
        H_{2m} - H_m = \ln 2 - \frac{1}{4m} +
        O\Bigl(\frac{1}{m^2}\Bigr),
        \qquad
        \sum_{j=1}^{m} \frac{1}{2j-1} = \frac{\ln m}{2} + \ln 2
        + \frac\gamma2 + o(1) ,
        \]
        dengan yang kedua lewat $\sum_{j \leq m} \frac{1}{2j-1} = H_{2m}
        - \frac12 H_m$, dan demikian pula $\sum_{j=1}^{m}
        \frac{1}{2j} = \frac{\ln m}{2} + \frac\gamma2 + o(1)$.
\end{enumerate}

\medskip
\noindent\textbf{Bagian III --- \omterm{def:b1:series:def}{Deret} harmonik berselang-seling,
sampai orde dua.}
\begin{enumerate}[resume]
  \item Tunjukkanlah (lewat induksi, atau pengelompokan) kesamaan
        $\sum_{k=1}^{2m} \frac{(-1)^{k-1}}{k} = H_{2m} - H_m$,
        lalu simpulkanlah baik jumlahnya $\ln 2$ (sekali lagi) maupun
        kecepatannya yang persis:
        \[
        \sum_{k=1}^{2m} \frac{(-1)^{k-1}}{k}
        = \ln 2 - \frac{1}{4m} + O\Bigl(\frac{1}{m^2}\Bigr) .
        \]
  \item Simpulkanlah galat asimtotik \omterm{def:b1:series:def}{deret} harmonik
        berselang-seling pada indeks \emph{mana pun}: $S - S_N \sim
        \frac{(-1)^N}{2N}$ --- yaitu dua kali lebih kecil daripada batas kasus
        terburuknya $a_{N+1} \approx \frac1N$ pada
        \cref{thm:b1:series:alternating}.
  \item (Percepatan yang cuma-cuma) Tunjukkanlah bahwa jumlah rata-ratanya
        $\tilde S_N = \frac{S_N + S_{N+1}}{2}$ memenuhi $\tilde
        S_N = \ln 2 + O\bigl(\frac{1}{N^2}\bigr)$. Periksalah: $S_{10}
        = 0.64563$, $S_{11} = 0.73654$, $\tilde S_{10} =
        0.69109$, terhadap $\ln 2 = 0.69315$: jadi satu perataan membeli
        dua tempat desimal.
  \item Jelaskanlah dalam dua kalimat mengapa tak ada muslihat semacam itu yang dapat membantu
        gejala berekor divergen yang \emph{positif} seperti apitan pertanyaan
        2: karena galat berselang-selingnya berayun (dengan tanda
        $(-1)^N$), sehingga perataannya meniadakan suku utamanya, sedangkan
        galat apitan $\gamma$-nya $\frac{1}{2n}$ bertanda tetap.
        (Adapun merata-ratakan $a_n$ dan $b_n$ \emph{memang} membantu:
        hubungkanlah $\frac{a_n + b_n}{2}$ dengan taksiran titik tengahnya
        $H_n - \ln\bigl(n + \frac12\bigr)$ lalu tunjukkanlah bahwa galatnya
        $O\bigl(\frac{1}{n^2}\bigr)$.)
\end{enumerate}

\medskip
\noindent\textbf{Bagian IV --- Kekakuan dan kegagalannya.}
\begin{enumerate}[resume]
  \item Tunjukkanlah bahwa bagian positif $\sum \frac{1}{2j-1}$ dan
        bagian negatif $\sum \frac{1}{2j}$ \omterm{def:b1:series:def}{deret} harmonik
        berselang-selingnya sama-sama divergen --- yaitu tanda tangan
        kekonvergenan yang \emph{bersyarat}.
  \item Buktikanlah \omterm{def:b1:logic:statement}{pernyataan} umum di balik pertanyaan 12: bahwa jika
        $\sum u_n$ konvergen tetapi $\sum \abs{u_n}$ divergen, maka
        \omterm{def:b1:series:def}{deret} bagian positifnya $\sum u_n^+$ dan bagian
        negatifnya $\sum u_n^-$ sama-sama divergen \emph{(dari
        $u_n^\pm = \frac{\abs{u_n} \pm u_n}{2}$: karena jika salah satunya
        konvergen, maka yang lain pun demikian, sehingga $\sum\abs{u_n}$)}.
        Adapun cadangan massa positif dan negatif yang tak habis-habis inilah
        yang akan dibelanjakan resep Riemann.
  \item (Kekakuan) Buktikanlah: bahwa jika $\sum u_n$ konvergen
        \emph{secara mutlak} dan $\sigma \colon \N \to \N$ sebuah
        bijeksi, maka $\sum u_{\sigma(n)}$ konvergen ke
        jumlah yang sama \emph{(karena untuk $N$ yang besar $M$ suku tersusun ulang
        yang pertama memuat $u_0, \dots, u_N$; lalu bandingkanlah jumlah parsialnya
        lewat ekornya $\sum_{n > N}\abs{u_n}$)}.
  \item (Resep Riemann) Misalkan $t \in \R$. Perikanlah penyusunan
        ulang yang rakus atas \omterm{def:b1:series:def}{deret} harmonik berselang-selingnya: ambillah
        suku positif $1, \frac13, \frac15, \dots$ sampai
        jumlah parsialnya pertama kali melampaui $t$, lalu suku negatif
        $-\frac12, -\frac14, \dots$ sampai ia pertama kali jatuh di bawah
        $t$, lalu ulangi. Tunjukkanlah bahwa setiap sukunya terpakai tepat
        sekali, bahwa setelah persilangan pertamanya jumlah parsialnya
        tinggal dalam jarak suku terakhir yang terpakai dari $t$, lalu simpulkanlah bahwa
        \omterm{def:b1:series:def}{deret} yang tersusun ulang itu konvergen ke $t$: sehingga \emph{sebarang}
        jumlah yang ditetapkan tercapai.
\end{enumerate}

\medskip
\noindent\textbf{Bagian V --- Rumus $(p, q)$-nya.} Tetapkanlah bilangan bulat
$p, q \geq 1$. Susunlah ulang \omterm{def:b1:series:def}{deret} harmonik berselang-selingnya dalam
blok: yaitu $p$ suku positif (yakni kebalikan bilangan ganjil berikutnya), lalu $q$
suku negatif (yakni kebalikan bilangan genap berikutnya), lalu ulangi.
\begin{enumerate}[resume]
  \item Periksalah bahwa ini penyusunan ulang yang sungguhan (yaitu setiap sukunya
        tepat sekali), dan bahwa untuk $(p, q) = (1, 2)$ ia berbunyi
        \[
        1 - \frac12 - \frac14 + \frac13 - \frac16 - \frac18 +
        \frac15 - \frac1{10} - \frac1{12} + \dots
        \]
  \item (Pemaruhan yang persis) Untuk $(p, q) = (1, 2)$, buktikanlah
        kesamaan bloknya
        \[
        \frac{1}{2k-1} - \frac{1}{4k-2} - \frac{1}{4k}
        = \frac12\Bigl(\frac{1}{2k-1} - \frac{1}{2k}\Bigr),
        \]
        lalu simpulkanlah hubungan yang \emph{persis} $T_{3K} = \frac12
        S_{2K}$ antara jumlah parsial yang tersusun ulang dan yang
        asli: sehingga pemaruhan jumlahnya kasatmata pada
        setiap tahap yang hingga, bukan hanya pada limitnya.
  \item Tunjukkanlah bahwa jumlah parsialnya setelah $K$ blok yang lengkap
        sama dengan $\sum_{j=1}^{pK} \frac{1}{2j-1} -
        \sum_{j=1}^{qK} \frac{1}{2j}$, lalu hitunglah limitnya
        dengan pertanyaan 7:
        \[
        \ln 2 + \frac12 \ln\frac pq .
        \]
  \item Kendalikanlah jumlah parsialnya \emph{di dalam} sebuah blok (karena
        sukunya menuju $0$) lalu simpulkanlah bahwa \omterm{def:b1:series:def}{deret} yang tersusun ulang
        secara $(p, q)$ konvergen ke $\ln 2 + \frac12\ln
        \frac pq$. Khususnya $(1, 2)$ memberikan $\frac{\ln
        2}{2}$: periksalah terhadap sembilan suku pertamanya, $T_9 =
        0.3083$, yang merayap menuju $0.3466$.
  \item Periksa kewarasan dan jangkauannya: bahwa $(1,1)$ memulihkan $\ln 2$;
        dan $(2,1)$ memberikan $\frac32\ln 2$; lalu jumlah mana yang tercapai
        lewat blok $(p, q)$, dan bagaimanakah menu terbilang ini
        dibandingkan dengan kartu penuh Riemann (pertanyaan 14)?
\end{enumerate}

\medskip
\noindent\textbf{Bagian VI --- Epilog: $\gamma$ dalam kerja, dan
sintesisnya.}
\begin{enumerate}[resume]
  \item Kenalilah jumlah \omterm{def:b1:series:def}{deret} yang konvergen
        $\sum_{k\geq1} \bigl(\frac1k - \ln\frac{k+1}{k}\bigr)$
        (\cref{exo:b1:series:7}): tunjukkanlah bahwa ia sama dengan $\gamma$.
  \item Jalankanlah resep Riemann (pertanyaan 14) bagi sasaran $t =
        1$ lalu daftarkanlah dua belas suku pertama yang dihasilkannya ($1, \frac13,
        -\frac12, \frac15, -\frac14, \frac17, \frac19, -\frac16,
        \frac1{11}, \frac1{13}, -\frac18, \frac1{15}$),
        sambil menghitung jumlah parsialnya ($\approx 0.980$) --- amatilah
        algoritmanya bernapas di sekitar sasarannya.
  \item Tunjukkanlah bahwa suatu \omterm{pb:b1:series:1}{penyusunan ulang deret} harmonik
        berselang-selingnya divergen ke $+\infty$ \emph{(lewat blok suku
        positif yang cukup panjang untuk memperoleh $1$ setiap kalinya, dengan memakai pertanyaan
        12, yang dipisahkan oleh suku negatif tunggal)}.
  \item Pertajamlah \cref{ex:b1:series:harmonicstack} dengan
        hukum $\gamma$-nya: tunjukkanlah bahwa indeks pertama dengan $H_N \geq
        20$ memenuhi $N = \eu^{\,20 - \gamma}\,(1 + o(1))
        \approx 2.7\cdot10^{8}$ --- jadi \omterm{pb:b1:series:1}{konstanta Euler} persis
        merupakan koreksi yang terlewat oleh apitan yang kasar itu.
  \item Sintesis, satu kalimat untuk masing-masing: (i) hukum $\gamma$-nya dan
        apa yang disumbangkan masing-masing dari ketiga kepingnya ($\ln n$, $\gamma$,
        $\frac{1}{2n}$); (ii) mengapakah kekonvergenan bersyarat
        membuat jumlahnya bergantung pada urutannya sedangkan \omterm{thm:b1:series:absolute}{kekonvergenan
        mutlak} melarangnya; (iii) bagaimanakah rumus $(p,q)$-nya
        merupakan sebuah \emph{perhitungan} dengan hukum $\gamma$-nya alih-alih
        klaim keberadaan yang abstrak; (iv) di manakah benang ini
        berlanjut --- yaitu \omterm{def:b1:series:def}{deret} pangkat dan hasil kali \omterm{def:b1:series:def}{deret}
        pada jilid Tahun ke-2, dan soal akhir pekan jilid Tahun ke-3
        tentang rumus Stirling, yang di situ pembukuan
        jumlah lawan \omterm{thm:b1:integration:def}{integral} yang sama berjalan dengan daya penuh.

\end{enumerate}
\end{problem}
